% Accelerated (anchored) version of the inexact primal-dual Douglas--Rachford
% method of Draft.tex, with proofs.
%
% This file compiles on its own (same preamble as Draft.tex).  The material
% between \begin{document} and \end{document} is written so that it can be
% pasted into Draft.tex as a section after the linear-system analysis.
%
% TODO (bibliography): the following entries are cited below but are not yet in
% bibliography.bib.  Verify each one before adding it.
%   halpern1967        B. Halpern, Fixed points of nonexpanding maps,
%                      Bull. Amer. Math. Soc. 73 (1967) 957--961.
%   lieder2021halpern  F. Lieder, On the convergence rate of the Halpern-iteration,
%                      Optim. Lett. 15 (2021) 405--418.
%   kim2021accelerated D. Kim, Accelerated proximal point method for maximally
%                      monotone operators, Math. Program. 190 (2021) 57--87.
%   park2022exact      J. Park and E. K. Ryu, Exact optimal accelerated complexity
%                      for fixed-point iterations, ICML 2022.
%   diakonikolas2020halpern  J. Diakonikolas, Halpern iteration for near-optimal
%                      and parameter-free monotone inclusion and strong solutions
%                      to variational inequalities, COLT 2020.
%   sabach2017first    S. Sabach and S. Shtern, A first order method for solving
%                      convex bilevel optimization problems, SIAM J. Optim. 27
%                      (2017) 640--660.
%   lu2024restarted    H. Lu and J. Yang, Restarted Halpern PDHG for linear
%                      programming, arXiv 2024.
%   robinson1981polyhedral  S. M. Robinson, Some continuity properties of
%                      polyhedral multifunctions, Math. Program. Study 14 (1981).

\documentclass[11pt]{article}

\usepackage{graphicx}
\usepackage{amssymb, amsfonts, amsmath, amsthm}
\usepackage[margin=1in, footskip=1cm, headheight = 1.5cm]{geometry}
\usepackage[sort]{natbib}
\usepackage{booktabs}
\usepackage{hyperref}
\hypersetup{
    colorlinks=true,
    linkcolor=blue,
    citecolor=cyan,
    pdftitle={Anchored acceleration of the inexact primal-dual DRS method},
    }
\bibliographystyle{unsrtnat}
\usepackage[nameinlink]{cleveref}

% Shared macros and environments for every document in tex/.
% Seeded from the group's standard template; project-specific additions are
% collected at the end.

% Commands
\newcommand{\eg}{{\it e.g.}}
\newcommand{\ie}{{\it i.e.}}

% Sets
\newcommand{\ones}{\mathbf 1}
\newcommand{\reals}{{\mathbf R}}
\newcommand{\integers}{{\mathbf  Z}}
\newcommand{\symm}{{\mathbf S}}  % symmetric matrices
\newcommand{\NPhard}{\mbox{$\mathcal{NP}$-hard}}
\newcommand{\prob}{{\mathbf P}}
\newcommand{\distrib}{{\mathcal P}}
\newcommand{\identity}{I}
\newcommand{\nullspace}{{\mathcal N}}
\newcommand{\range}{{\mathcal R}}
\newcommand{\Rank}{\mathop{\bf rank}}
\newcommand{\Tr}{\mathop{\bf Tr}}
\newcommand{\diag}{\mathop{\bf diag}}
\newcommand{\lambdamax}{{\lambda_{\rm max}}}
\newcommand{\lambdamin}{\lambda_{\rm min}}
\newcommand{\Expect}{\mathop{\bf E{}}}
\newcommand{\Prob}{\mathop{\bf prob}}
\newcommand{\Co}{{\mathop {\bf Co}}} % convex hull
\newcommand{\dist}{\mathop{\bf dist{}}}
\newcommand{\argmin}{\mathop{\rm argmin}}
\newcommand{\argmax}{\mathop{\rm argmax}}
\newcommand{\epi}{\mathop{\bf epi}} % epigraph
\newcommand{\Vol}{\mathop{\bf vol}}
\newcommand{\dom}{\mathop{\bf dom}} % domain
\newcommand{\intr}{\mathop{\bf int}}

\newcommand{\sign}{\mathop{\bf sign}}
\newcommand{\round}{\mathop{\bf round}}

% Sections (for citation)
\newcommand{\Sec}{Section\;}

% Create theorems and other environments
\newtheorem{theorem}{Theorem}[section]  % Restart counter every section
\newtheorem{lemma}{Lemma}[section]  % Restart counter every section
\newtheorem{corollary}{Corollary}[theorem]  % Restart counter every theorem
\newtheorem{proposition}{Proposition}[section]
\newtheorem{assumption}{Assumption}[section]

\theoremstyle{definition}
\newtheorem{definition}{Definition}[section]
\newtheorem{problem}{Problem}[section]

\theoremstyle{remark}
\newtheorem{example}{Example}[section]
\newtheorem{exercise}{Exercise}[section]
\newtheorem{remark}{Remark}[section]

% Acronyms.  Defined only when the preamble loaded the glossaries package, so
% that this file can be \input from a document that does not use acronyms.
\makeatletter
\@ifundefined{newacronym}{}{%
  \newacronym{LP}{LP}{linear program}
  \newacronym{QP}{QP}{quadratic program}
  \newacronym{MIQP}{MIQP}{mixed-integer quadratic program}
  \newacronym{MIP}{MIP}{mixed-integer program}
  \newacronym{MILP}{MILP}{mixed-integer linear program}
  \newacronym{MINLP}{MINLP}{mixed-integer nonlinear program}
  \newacronym{sBB}{sBB}{spatial branch and bound}
  \newacronym{PWA}{PWA}{piecewise affine}
  \newacronym{SVM}{SVM}{support vector machines}
  \newacronym{ReLU}{ReLU}{rectified linear unit}
  \newacronym{CPU}{CPU}{central processing unit}
  \newacronym{GPU}{GPU}{graphics processing unit}
  \newacronym{MPC}{MPC}{model predictive control}
  \newacronym{ADMM}{ADMM}{alternating direction method of multipliers}
  \newacronym{ADP}{ADP}{approximate dynamic programming}
  \newacronym{FPGA}{FPGA}{field-programmable gate array}
  \newacronym{DRS}{DRS}{Douglas--Rachford splitting}
  \newacronym{RandNLA}{RandNLA}{randomized numerical linear algebra}
  \newacronym{CG}{CG}{conjugate gradients}
  \newacronym{GMRES}{GMRES}{generalized minimal residual method}
}
\makeatother

% Project-specific definitions
\newcommand{\prox}{\mathrm{prox}}
\newcommand{\Fix}{\mathrm{Fix}}
\newcommand{\R}{\mathbb{R}}


\crefname{equation}{}{}
\crefname{assumption}{assumption}{assumptions}

\usepackage{xcolor}
\newcommand{\mynote}[1]{\textcolor{blue}{[#1]}}

\title{An Anchored Acceleration of the Inexact Primal-Dual Douglas--Rachford Method}
\author{Pranav Reddy}
\date{\small \today \vspace{-5mm}}

\begin{document}

\maketitle

\begin{abstract}
    We accelerate the inexact primal-dual Douglas--Rachford method of the main draft by anchoring it to its starting point in the style of the Halpern iteration.
    The anchored method uses the same inexact linear solve, the same exact proximal step, and the same relative-error acceptance test.
    With exact linear solves, the fixed-point residual satisfies $\|s^{(k+1)}\|_M \le 2\|z^{(0)} - z^\star\|_M/(\theta k + 2)$, which is the best possible rate for this class of methods.
    With inexact solves whose relative tolerances are square summable, the residual still decays as $O(1/k)$, compared to $O(1/\sqrt{k})$ for the base method.
    Restarting the anchor whenever the residual has decreased by a fixed factor gives linear convergence under an error-bound condition, with a number of outer iterations proportional to the error-bound constant rather than to its square.
\end{abstract}

\section{Overview}\label{sec:acc-overview}
The main draft studies the relaxed iteration $z^{(k+1)} = z^{(k)} + \theta s^{(k+1)}$, where $s^{(k+1)}$ is the (inexactly computed) Douglas--Rachford step at $z^{(k)}$.
With exact linear solves this is the Krasnosel'ski\u{\i}--Mann iteration of a firmly nonexpansive map, and the telescoping bound in the draft gives $\min_{t \le k}\|s^{(t+1)}\|_M = O(1/\sqrt{k})$.
For nonexpansive maps, the Halpern iteration \citep{halpern1967} improves this rate to $O(1/k)$ \citep{sabach2017first, lieder2021halpern, diakonikolas2020halpern}, and the rate $O(1/k)$ is optimal among all methods that only evaluate the map \citep{park2022exact}.
The version for resolvents of maximally monotone operators is the accelerated proximal point method of \citet{kim2021accelerated}.

We apply the same idea to the inexact method.
The only change is the outer update, which becomes
\begin{align}
    z^{(k+1)} = \frac{1}{k+2}\, z^{(0)} + \frac{k+1}{k+2}\left(z^{(k)} + \theta s^{(k+1)}\right), \label{eq:anchored-update}
\end{align}
so that every iterate is pulled back toward the starting point with a weight that decays like $1/k$.
The inner solve, the proximal step, and the acceptance test $\|\varepsilon^{(k+1)}\|_M \le \sigma_k \|s^{(k+1)}\|_M$ are unchanged, except that the tolerance $\sigma_k$ may now depend on $k$.
\Cref{sec:acc-oracle} isolates the two properties of the inexact step that the analysis uses.
\Cref{sec:acc-anchored} proves the $O(1/k)$ rate.
\Cref{sec:acc-restart} adds restarts and proves linear convergence under an error bound.

\section{The inexact step as a monotone oracle}\label{sec:acc-oracle}
Throughout, $\langle \cdot, \cdot \rangle_M$ and $\|\cdot\|_M$ denote the inner product and norm induced by $M = \diag(\gamma_x^{-1} I, \gamma_\lambda^{-1} I)$, and $z = (x, \lambda) \in \R^{n+m}$.
Fix $z = (x,\lambda)$ and let $(\xi, \nu)$ be any approximate solution of the linear system
\begin{align}
    \begin{bmatrix} I & \gamma_x A^T \\ -\gamma_\lambda A & I \end{bmatrix}
    \begin{bmatrix} \xi \\ \nu \end{bmatrix}
    = \begin{bmatrix} x \\ \lambda - \gamma_\lambda b \end{bmatrix} \label{eq:acc-linear-system}
\end{align}
at $z$, which is the system solved at every outer iteration of the draft.
As in the draft, define
\begin{align}
    v &= \prox_{\gamma_x f}(\xi - \gamma_x A^T \nu), \qquad
    \hat{z} = \begin{bmatrix} v + \gamma_x A^T \nu \\ \nu \end{bmatrix}, \qquad
    s = \begin{bmatrix} v - \xi \\ \gamma_\lambda (A\xi - b) \end{bmatrix}, \label{eq:acc-step-definitions}
\end{align}
and let $\varepsilon = \hat{z} - z - s$ be the residual of the linear system.
We call $(\hat z, s)$ an \emph{admissible pair} if it arises from some $(\xi, \nu) \in \R^n \times \R^m$ through \Cref{eq:acc-step-definitions}, and we write $\mathcal{G}$ for the set of admissible pairs.
Note that $\mathcal{G}$ does not depend on $z$: the point $z$ only enters through the residual $\varepsilon$.
Let $\mathcal{Z}^\star = \{ (x^\star + \gamma_x A^T\lambda^\star, \lambda^\star) \mid (x^\star,\lambda^\star) \text{ satisfies } -A^T\lambda^\star \in \partial f(x^\star),\ Ax^\star = b \}$ denote the set of fixed points.
Finally, let $s^\star(z)$ denote the step \Cref{eq:acc-step-definitions} obtained from the exact solution of the linear system at $z$; this is the step of the exact Douglas--Rachford method, and $s^\star(z) = 0$ if and only if $z \in \mathcal{Z}^\star$.

\begin{lemma}\label{lem:acc-monotone}
    The following statements hold.
    \begin{enumerate}
        \item For any two admissible pairs $(\hat z_1, s_1)$ and $(\hat z_2, s_2)$,
        \begin{align}
            \langle \hat z_1 - \hat z_2, s_1 - s_2\rangle_M \le 0. \label{eq:acc-two-point}
        \end{align}
        \item For every $z^\star \in \mathcal{Z}^\star$, the pair $(z^\star, 0)$ is admissible.
        \item Let $(\hat z, s)$ be admissible with residual $\varepsilon$ at $z$, and let $(\hat z^\star, s^\star(z))$ be the exact pair at $z$.
        Then $\|s - s^\star(z)\|_M \le \|\varepsilon\|_M$.
        In particular, if $\|\varepsilon\|_M \le \sigma\|s\|_M$ with $\sigma < 1$, then
        \begin{align}
            (1-\sigma)\|s\|_M \le \|s^\star(z)\|_M \le (1+\sigma)\|s\|_M. \label{eq:acc-exact-vs-inexact}
        \end{align}
    \end{enumerate}
\end{lemma}
\begin{proof}
    For an admissible pair generated by $(\xi_i, \nu_i)$, the optimality condition of the proximal step reads
    \begin{align}
        u_i := \gamma_x^{-1}(\xi_i - v_i) - A^T\nu_i \in \partial f(v_i). \label{eq:acc-prox-optimality}
    \end{align}
    Expanding the $M$-inner product and cancelling the terms in $A\xi_i$,
    \begin{align}
        \langle \hat z_1 - \hat z_2, s_1 - s_2\rangle_M
        &= \gamma_x^{-1}\langle v_1 - v_2 + \gamma_x A^T(\nu_1 - \nu_2),\, (v_1 - \xi_1) - (v_2 - \xi_2)\rangle
           + \langle \nu_1 - \nu_2,\, A(\xi_1 - \xi_2)\rangle \\
        &= \gamma_x^{-1}\langle v_1 - v_2,\, (v_1 - \xi_1) - (v_2 - \xi_2)\rangle
           + \langle \nu_1 - \nu_2,\, A(v_1 - v_2)\rangle \\
        &= \left\langle v_1 - v_2,\, \gamma_x^{-1}\big((v_1 - \xi_1) - (v_2 - \xi_2)\big) + A^T(\nu_1 - \nu_2)\right\rangle \\
        &= -\langle v_1 - v_2,\, u_1 - u_2\rangle \le 0,
    \end{align}
    where the last inequality is monotonicity of $\partial f$.
    This proves the first statement.

    For the second statement take $\xi = x^\star$ and $\nu = \lambda^\star$.
    Since $-A^T\lambda^\star \in \partial f(x^\star)$, the proximal step returns $v = \prox_{\gamma_x f}(x^\star - \gamma_x A^T\lambda^\star) = x^\star$.
    Hence $\hat z = (x^\star + \gamma_x A^T\lambda^\star, \lambda^\star) = z^\star$ and $s = (0, \gamma_\lambda(Ax^\star - b)) = 0$.

    For the third statement, $\hat z = z + s + \varepsilon$ and $\hat z^\star = z + s^\star(z)$ give $\varepsilon = (\hat z - \hat z^\star) - (s - s^\star(z))$.
    Expanding the square and using \Cref{eq:acc-two-point},
    \begin{align}
        \|\varepsilon\|_M^2 = \|\hat z - \hat z^\star\|_M^2 - 2\langle \hat z - \hat z^\star, s - s^\star(z)\rangle_M + \|s - s^\star(z)\|_M^2 \ge \|s - s^\star(z)\|_M^2.
    \end{align}
    The two-sided bound \Cref{eq:acc-exact-vs-inexact} follows from the triangle inequality.
\end{proof}

The first statement says that $\mathcal{G}$ is the graph of a monotone operator (up to the sign of $s$); it is the operator whose resolvent is the exact Douglas--Rachford map, as in \citet{EcksteinRelativeError2018}.
With $z^\star \in \mathcal{Z}^\star$ and $\hat z = z + s + \varepsilon$, inequality \Cref{eq:acc-two-point} gives
\begin{align}
    \langle s, z - z^\star\rangle_M \le -\|s\|_M^2 - \langle \varepsilon, s\rangle_M \le -(1-\sigma)\|s\|_M^2 \label{eq:acc-star-inequality}
\end{align}
whenever $\|\varepsilon\|_M \le \sigma\|s\|_M$.
This is the only place where the acceptance test enters the exact analysis of the draft; the anchored analysis below also uses \Cref{eq:acc-two-point} between two consecutive iterations.

\section{The anchored method}\label{sec:acc-anchored}

\subsection{Algorithm}
Let $\theta \in (0,2]$, let $(\sigma_k)_{k \ge 0}$ be tolerances in $[0,1)$, and let $z^{(0)} = (x^{(0)}, \lambda^{(0)})$.
For $k = 0, 1, 2, \ldots$:
\begin{enumerate}
    \item Solve the linear system \Cref{eq:acc-linear-system} at $z^{(k)}$ inexactly, obtaining $(\xi^{(k+1)}, \nu^{(k+1)})$, and form $v^{(k+1)}$, $\hat z^{(k+1)}$, $s^{(k+1)}$, and $\varepsilon^{(k+1)}$ as in \Cref{eq:acc-step-definitions}.
    Continue the inner solver until $\|\varepsilon^{(k+1)}\|_M \le \sigma_k\|s^{(k+1)}\|_M$.
    \item Update
    \begin{subequations}\label{eq:acc-anchored-xy}
    \begin{align}
        x^{(k+1)} &= \frac{1}{k+2}\, x^{(0)} + \frac{k+1}{k+2}\left(x^{(k)} + \theta\,(v^{(k+1)} - \xi^{(k+1)})\right), \\
        \lambda^{(k+1)} &= \frac{1}{k+2}\, \lambda^{(0)} + \frac{k+1}{k+2}\left(\lambda^{(k)} - \theta\gamma_\lambda\,(b - A\xi^{(k+1)})\right).
    \end{align}
    \end{subequations}
\end{enumerate}
In the notation $z^{(k)} = (x^{(k)}, \lambda^{(k)})$ this is exactly \Cref{eq:anchored-update}.
The per-iteration cost is that of the base method plus two vector additions.
With exact solves ($\sigma_k \equiv 0$) and $\theta = 1$ the method is the Halpern iteration of the exact Douglas--Rachford map; with $\theta = 2$ it is the Halpern iteration of its reflection, which is the accelerated proximal point method of \citet{kim2021accelerated}.

\subsection{A potential function}
Write $D = \|z^{(0)} - z^\star\|_M$ for a fixed $z^\star \in \mathcal{Z}^\star$; since $z^\star$ is arbitrary, all bounds below hold with $D = \dist_M(z^{(0)}, \mathcal{Z}^\star)$.
Define
\begin{align}
    V_k = -(k+1)\left\langle s^{(k+1)}, z^{(k)} - z^{(0)}\right\rangle_M + \frac{\theta}{2}\, k(k+1)\, \|s^{(k+1)}\|_M^2. \label{eq:acc-potential}
\end{align}
The following one-step inequality is the core of the analysis.

\begin{lemma}\label{lem:acc-one-step}
    Let $\Lambda_k = (k+1)(k+2)$, $d_k = s^{(k+1)} - s^{(k+2)}$, and $e_k = \varepsilon^{(k+1)} - \varepsilon^{(k+2)}$.
    For every $k \ge 0$,
    \begin{align}
        V_{k+1} \le V_k - \frac{2-\theta}{2}\,\Lambda_k\,\|d_k\|_M^2 - \Lambda_k\,\langle e_k, d_k\rangle_M. \label{eq:acc-one-step}
    \end{align}
\end{lemma}
\begin{proof}
    Write $r_k = z^{(k)} - z^{(0)}$, $\beta_k = 1/(k+2)$, $a = \|s^{(k+1)}\|_M^2$, $b = \|s^{(k+2)}\|_M^2$, and $p = \langle s^{(k+1)}, s^{(k+2)}\rangle_M$.
    The update \Cref{eq:anchored-update} gives the two identities
    \begin{align}
        (k+2)\, r_{k+1} &= (k+1)\left(r_k + \theta s^{(k+1)}\right), \label{eq:acc-anchor-identity}\\
        z^{(k)} - z^{(k+1)} &= \beta_k r_k - (1-\beta_k)\,\theta s^{(k+1)}. \label{eq:acc-difference-identity}
    \end{align}
    Since $\hat z^{(k+1)} = z^{(k)} + s^{(k+1)} + \varepsilon^{(k+1)}$ and $\hat z^{(k+2)} = z^{(k+1)} + s^{(k+2)} + \varepsilon^{(k+2)}$, \Cref{eq:acc-difference-identity} gives
    \begin{align}
        \hat z^{(k+1)} - \hat z^{(k+2)} = \beta_k r_k - (1-\beta_k)\theta s^{(k+1)} + d_k + e_k.
    \end{align}
    Apply \Cref{eq:acc-two-point} to the pairs $(\hat z^{(k+1)}, s^{(k+1)})$ and $(\hat z^{(k+2)}, s^{(k+2)})$ and multiply by $\Lambda_k > 0$.
    Using $\beta_k\Lambda_k = k+1$ and $(1-\beta_k)\Lambda_k = (k+1)^2$,
    \begin{align}
        (k+1)\langle r_k, s^{(k+1)}\rangle_M - (k+1)\langle r_k, s^{(k+2)}\rangle_M - \theta(k+1)^2(a - p) \qquad & \notag \\
        {} + \Lambda_k\|d_k\|_M^2 + \Lambda_k\langle e_k, d_k\rangle_M &\le 0. \label{eq:acc-scaled-monotone}
    \end{align}
    By \Cref{eq:acc-anchor-identity},
    \begin{align}
        V_{k+1} = -(k+1)\langle r_k, s^{(k+2)}\rangle_M - \theta(k+1)\,p + \frac{\theta}{2}\Lambda_k\, b.
    \end{align}
    Bounding the first term with \Cref{eq:acc-scaled-monotone} and using $\frac{\theta}{2}k(k+1) = \theta(k+1)^2 - \frac{\theta}{2}\Lambda_k$,
    \begin{align}
        V_{k+1}
        &\le -(k+1)\langle r_k, s^{(k+1)}\rangle_M + \theta(k+1)^2 a - \theta(k+1)^2 p - \theta(k+1)p + \frac{\theta}{2}\Lambda_k b \notag \\
        &\qquad - \Lambda_k\|d_k\|_M^2 - \Lambda_k\langle e_k, d_k\rangle_M \\
        &= V_k + \frac{\theta}{2}\Lambda_k\,(a - 2p + b) - \Lambda_k\|d_k\|_M^2 - \Lambda_k\langle e_k, d_k\rangle_M,
    \end{align}
    where we used $\theta(k+1)^2 + \theta(k+1) = \theta\Lambda_k$.
    Since $a - 2p + b = \|d_k\|_M^2$, this is \Cref{eq:acc-one-step}.
\end{proof}

The lower bound on $V_k$ comes from \Cref{eq:acc-star-inequality}.
Writing $z^{(k)} - z^{(0)} = (z^{(k)} - z^\star) + (z^\star - z^{(0)})$ and using Cauchy--Schwarz on the second term,
\begin{align}
    V_k \ge W_k\, \|s^{(k+1)}\|_M^2 - (k+1)\, D\, \|s^{(k+1)}\|_M,
    \qquad W_k := (k+1)(1 - \sigma_k) + \frac{\theta}{2}k(k+1). \label{eq:acc-potential-lower}
\end{align}

\subsection{Exact linear solves}
\begin{theorem}\label{thm:acc-exact}
    Let $\theta \in (0,2]$ and $\sigma_k \equiv 0$.
    Then for every $k \ge 0$,
    \begin{align}
        \|s^{(k+1)}\|_M \le \frac{2D}{\theta k + 2}.
    \end{align}
\end{theorem}
\begin{proof}
    With $e_k = 0$ and $\theta \le 2$, \Cref{eq:acc-one-step} gives $V_{k+1} \le V_k$, hence $V_k \le V_0 = 0$ because $z^{(0)} - z^{(0)} = 0$.
    Combining with \Cref{eq:acc-potential-lower} and dividing by $(k+1)\|s^{(k+1)}\|_M$,
    \begin{align}
        \left(1 + \frac{\theta k}{2}\right)\|s^{(k+1)}\|_M \le D.
    \end{align}
\end{proof}

For $\theta = 1$ the bound reads $\|s^{(k+1)}\|_M \le 2D/(k+2)$, which is the bound of \citet{lieder2021halpern} for the Halpern iteration of a nonexpansive map, and for $\theta = 2$ it reads $D/(k+1)$, the bound of \citet{kim2021accelerated} for the accelerated proximal point method.
Both are known to be tight \citep{park2022exact}, so \Cref{thm:acc-exact} cannot be improved by a different proof.
\mynote{TODO: I also solved the worst-case (performance estimation) problem for this iteration numerically for $\theta \in \{1, 2\}$ and $k \le 8$; the computed worst cases match $2D/(\theta k + 2)$ to four digits.  The script lives only in scratch space; add it under \texttt{code/experiments/} if we keep this remark.}

\subsection{Inexact linear solves}
\Cref{lem:acc-one-step} shows exactly how the inexactness enters: the error term $\Lambda_k\langle e_k, d_k\rangle_M$ is of the same order $k^2\|s\|_M^2$ as the potential itself, while the slack $\frac{2-\theta}{2}\Lambda_k\|d_k\|_M^2$ only controls the \emph{difference} of consecutive steps.
A constant tolerance therefore cannot be absorbed by this argument (see \Cref{rem:acc-fixed-sigma}), but a square-summable tolerance sequence can.

\begin{theorem}\label{thm:acc-inexact}
    Let $\theta \in (0,2)$ and define
    \begin{align}
        \kappa_\theta = \max\{4,\, 2/\theta\}, \qquad
        c_\theta = \frac{2\kappa_\theta}{2-\theta}, \qquad
        m_\theta = \max\{1,\, 1/\theta\}.
    \end{align}
    Suppose the tolerances satisfy $c_\theta\sigma_k^2 \le 1/2$ for all $k$ and $\sum_{k\ge 0}\sigma_k^2 \le \bar\Sigma$.
    Then for every $k \ge 0$,
    \begin{align}
        \|s^{(k+1)}\|_M \le \frac{2D}{\theta k + 1} + \sqrt{2 m_\theta}\; e^{3 c_\theta\bar\Sigma/2}\, \frac{D}{\sqrt{(k+1)(\theta k + 1)}}. \label{eq:acc-inexact-rate}
    \end{align}
\end{theorem}
\begin{proof}
    Write $a_k = \|s^{(k+1)}\|_M^2$.
    The condition $c_\theta\sigma_k^2 \le 1/2$ implies $\sigma_k \le 1/2$, since $c_\theta \ge 8/(2-\theta) \ge 4$.
    Hence $W_k \ge (k+1)(\theta k + 1)/2$ in \Cref{eq:acc-potential-lower}.

    \emph{Step 1: error term.}
    By Young's inequality, $-\langle e_k, d_k\rangle_M \le \frac{2-\theta}{2}\|d_k\|_M^2 + \frac{1}{2(2-\theta)}\|e_k\|_M^2$, so \Cref{eq:acc-one-step} gives $V_{k+1} \le V_k + \Lambda_k\|e_k\|_M^2/(2(2-\theta))$.
    The acceptance test gives $\|e_k\|_M \le \sigma_k\sqrt{a_k} + \sigma_{k+1}\sqrt{a_{k+1}}$, and therefore
    \begin{align}
        V_{k+1} \le V_k + \frac{\Lambda_k}{2-\theta}\left(\sigma_k^2 a_k + \sigma_{k+1}^2 a_{k+1}\right). \label{eq:acc-step1}
    \end{align}

    \emph{Step 2: a nonnegative companion of $V_k$.}
    Define $P_k = -(k+1)\langle s^{(k+1)}, z^{(k)} - z^\star\rangle_M + \frac{\theta}{2}k(k+1)a_k$.
    By \Cref{eq:acc-star-inequality}, $P_k \ge W_k a_k \ge 0$, and $V_k = P_k - (k+1)\langle s^{(k+1)}, z^\star - z^{(0)}\rangle_M$.
    Since $(k+1)^2/W_k \le 2(k+1)/(\theta k+1) \le 2m_\theta$, Young's inequality gives
    \begin{align}
        (k+1) D \sqrt{a_k} \le \frac{W_k a_k}{2} + \frac{(k+1)^2 D^2}{2 W_k} \le \frac{P_k}{2} + m_\theta D^2,
    \end{align}
    so $|V_k - P_k| \le P_k/2 + m_\theta D^2$.
    Let $\tilde V_k = V_k + m_\theta D^2$; then $0 \le P_k/2 \le \tilde V_k$.
    Next, $\Lambda_k/W_k \le 2(k+2)/(\theta k + 1) \le \kappa_\theta$, because $k \mapsto 2(k+2)/(\theta k+1)$ is monotone with values $4$ at $k=0$ and $2/\theta$ at infinity, and $\Lambda_k/W_{k+1} \le 2(k+1)/(\theta(k+1)+1) \le 2/\theta \le \kappa_\theta$.
    Together with $a_k \le P_k/W_k \le 2\tilde V_k/W_k$ this turns \Cref{eq:acc-step1} into
    \begin{align}
        \tilde V_{k+1} \le \tilde V_k + c_\theta\left(\sigma_k^2 \tilde V_k + \sigma_{k+1}^2 \tilde V_{k+1}\right). \label{eq:acc-step2}
    \end{align}

    \emph{Step 3: the potential stays bounded.}
    Rearranging \Cref{eq:acc-step2} and using $1 + t \le e^t$ and $(1-t)^{-1} \le e^{2t}$ for $t \in [0, 1/2]$,
    \begin{align}
        \tilde V_{k+1} \le \frac{1 + c_\theta\sigma_k^2}{1 - c_\theta\sigma_{k+1}^2}\, \tilde V_k \le e^{c_\theta\sigma_k^2 + 2c_\theta\sigma_{k+1}^2}\, \tilde V_k .
    \end{align}
    Multiplying out and using $\tilde V_0 = m_\theta D^2$ (because $V_0 = 0$),
    \begin{align}
        \tilde V_k \le e^{3 c_\theta \bar\Sigma}\, m_\theta D^2 =: E \qquad \text{for all } k \ge 0. \label{eq:acc-step3}
    \end{align}

    \emph{Step 4: conclusion.}
    By \Cref{eq:acc-potential-lower} and \Cref{eq:acc-step3}, $W_k a_k - (k+1)D\sqrt{a_k} \le V_k \le E$.
    Solving this quadratic inequality in $\sqrt{a_k}$ and using $\sqrt{u + w} \le \sqrt u + \sqrt w$,
    \begin{align}
        \sqrt{a_k} \le \frac{(k+1)D + \sqrt{(k+1)^2D^2 + 4W_kE}}{2W_k} \le \frac{(k+1)D}{W_k} + \sqrt{\frac{E}{W_k}}.
    \end{align}
    Inserting $W_k \ge (k+1)(\theta k+1)/2$ gives \Cref{eq:acc-inexact-rate}.
\end{proof}

\begin{corollary}\label{cor:acc-kkt}
    Under the assumptions of \Cref{thm:acc-inexact}, let $C_k$ denote the right-hand side of \Cref{eq:acc-inexact-rate}.
    The primal-dual pair $(v^{(k+1)}, \nu^{(k+1)})$ satisfies
    \begin{align}
        \dist\!\left(0,\ \partial f(v^{(k+1)}) + A^T\nu^{(k+1)}\right) \le \gamma_x^{-1/2}\, C_k,
        \qquad
        \|A v^{(k+1)} - b\| \le \left(\|A\|\gamma_x^{1/2} + \gamma_\lambda^{-1/2}\right) C_k .
    \end{align}
    In particular both KKT residuals decay as $O(1/k)$.
\end{corollary}
\begin{proof}
    By \Cref{eq:acc-prox-optimality}, $-\gamma_x^{-1}(v^{(k+1)} - \xi^{(k+1)}) \in \partial f(v^{(k+1)}) + A^T\nu^{(k+1)}$.
    Since $\|s\|_M^2 = \gamma_x^{-1}\|v - \xi\|^2 + \gamma_\lambda^{-1}\|\gamma_\lambda(A\xi - b)\|^2$, we have $\|v - \xi\| \le \gamma_x^{1/2}\|s\|_M$ and $\|A\xi - b\| \le \gamma_\lambda^{-1/2}\|s\|_M$.
    The primal bound follows from $\|Av - b\| \le \|A\|\|v - \xi\| + \|A\xi - b\|$.
\end{proof}

\begin{remark}[Choice of tolerances and inner work]\label{rem:acc-inner-work}
    A convenient schedule is $\sigma_k = \bar\sigma\,(k+1)^{-(1/2+\delta)}$ with $\delta > 0$ and $c_\theta\bar\sigma^2 \le 1/2$, for which $\bar\Sigma = \bar\sigma^2\sum_{k\ge 1}k^{-(1+2\delta)} < \infty$.
    The stopping-time bound of the draft applies verbatim with $\sigma$ replaced by $\sigma_k$, and $\beta_k$ scales as $\sigma_k^{-2}$.
    Hence the expected number of inner iterations at outer iteration $k$ grows only like $(1+2\delta)\log(k+1)/\log(1/a)$ plus a constant, while the outer residual decays like $1/k$.
    To reach $\|s\|_M \le \epsilon$ the base method needs on the order of $\epsilon^{-2}$ outer iterations with $O(1)$ inner iterations each, whereas the anchored method needs on the order of $\epsilon^{-1}$ outer iterations with $O(\log(1/\epsilon))$ inner iterations each.
    \mynote{TODO: the warm-start lemma of the draft bounds $\mathcal{E}_{k+1,0}$ in terms of $\|s^{(k)}\|_M$ using the update $z^{(k+1)} - z^{(k)} = \theta s^{(k+1)}$.  For the anchored update the drift is $z^{(k+1)} - z^{(k)} = \frac{k+1}{k+2}\theta s^{(k+1)} - \frac{1}{k+2}(z^{(k)} - z^{(0)})$, so that lemma needs to be redone before the total-work theorem can be transferred.}
\end{remark}

\begin{remark}[Constant tolerance]\label{rem:acc-fixed-sigma}
    For a constant tolerance $\sigma_k \equiv \sigma > 0$, \Cref{eq:acc-step1} only gives $V_k \le O(\sigma^2 D^2 k)$ when $\|s^{(k+1)}\|_M = O(D/k)$, which through \Cref{eq:acc-potential-lower} yields no more than the $O(1/\sqrt{k})$ rate of the base method.
    This is not an artifact of a careless estimate: for every anchor sequence of the form $1 - \beta_k = \tau_k/\tau_{k+1}$ and every potential of the form $\tau_k\langle s^{(k+1)}, z^{(k)} - z^{(0)}\rangle_M + A_k\|s^{(k+1)}\|_M^2$, the coefficients that make the exact potential nonincreasing are determined at leading order, so there is no room of order $k^2\|s\|_M^2$ to absorb an error of that order.
    Whether the anchored method with constant $\sigma$ still converges at rate $O(1/k)$ is open.
    \mynote{TODO: numerical worst-case (performance estimation) computations for $\theta=1$, over all firmly nonexpansive maps and all residuals with $\|\varepsilon\|_M \le \sigma\|s\|_M$, gave the following values of $k\|s^{(k)}\|_M/D$ (nonconvex Burer--Monteiro solve, so lower bounds on the true worst case; the $\sigma = 0$ column matches $2k/(k+1)$ exactly).  Reproduce with a proper SDP solver before quoting.}
    \begin{center}
    \begin{tabular}{lccccc}
        \toprule
        $k$ & 1 & 2 & 4 & 8 & 12 \\
        \midrule
        $\sigma = 0$   & 1.00 & 1.33 & 1.60 & 1.78 & 1.85 \\
        $\sigma = 0.1$ & 1.11 & 1.44 & 1.69 & 1.87 & 1.94 \\
        $\sigma = 0.3$ & 1.43 & 1.82 & 2.21 & 2.52 & 2.67 \\
        \bottomrule
    \end{tabular}
    \end{center}
    The growth with $k$ for $\sigma = 0.3$ is slow but has not clearly saturated by $k = 12$.
\end{remark}

\section{Restarts and linear convergence}\label{sec:acc-restart}
The anchored method converges at rate $O(1/k)$ in the worst case, but it also converges no faster than that on easy problems.
Indeed, by \Cref{eq:acc-anchor-identity}, $(k+1)(z^{(k)} - z^{(0)}) = \theta\sum_{j=1}^{k} j\, s^{(j)}$, so if the iterates converge to some $z^{(\infty)}$ then $\sum_{j \le k} j\, s^{(j)} \approx (k+1)(z^{(\infty)} - z^{(0)})/\theta$, which forces $\|s^{(k)}\|_M$ to be of order $\|z^{(\infty)} - z^{(0)}\|_M/(\theta k)$.
The base method, in contrast, converges linearly on linear and quadratic programs.
The standard remedy is to restart the anchor whenever the residual has dropped by a fixed factor, as is done for the primal-dual hybrid gradient method in \citet{lu2024restarted}.
We now show that, under an error-bound condition, the restarted method converges linearly, and we compare the resulting iteration count with that of the base method.

\begin{assumption}[Error bound]\label{ass:acc-error-bound}
    There is $\eta > 0$ such that $\dist_M(z, \mathcal{Z}^\star) \le \eta\,\|s^\star(z)\|_M$ for all $z \in \R^{n+m}$.
\end{assumption}
For linear programs, and more generally when $f$ is piecewise linear-quadratic, the exact Douglas--Rachford map is piecewise affine, and \Cref{ass:acc-error-bound} holds by the theory of polyhedral multifunctions \citep{robinson1981polyhedral}.
\mynote{TODO: verify this citation and the exact statement; a local version of the assumption on a level set containing the iterates suffices for the theorem below, with the obvious changes.}

\subsection{Algorithm}
Fix $\rho \in (0,1)$, $\theta \in (0,2)$, and a nonincreasing tolerance sequence $(\sigma_j)_{j\ge 0}$ satisfying the hypotheses of \Cref{thm:acc-inexact}.
The method runs the anchored iteration in epochs.
An epoch starts from an anchor $\bar z$ with a step $\bar s$ that was computed at $\bar z$ with tolerance at most $\sigma_0$, and sets $z^{(0)} = \bar z$, $s^{(1)} = \bar s$, and $j = 0$.
It then repeats:
\begin{enumerate}
    \item Form $z^{(j+1)}$ from $z^{(j)}$ and $s^{(j+1)}$ by \Cref{eq:anchored-update} with $k = j$, and increase $j$ by one.
    \item Compute the inexact step $s^{(j+1)}$ at $z^{(j)}$ with tolerance $\sigma_j$.
    \item If $\|s^{(j+1)}\|_M \le \rho\,\|s^{(1)}\|_M$, start a new epoch with anchor $\bar z = z^{(j)}$ and step $\bar s = s^{(j+1)}$.
\end{enumerate}
The first epoch is initialized with $\bar z$ equal to the starting point and $\bar s$ the step computed there with tolerance $\sigma_0$.
Each linear solve is used exactly once, so the cost per outer iteration is the same as for the anchored method.
Because $(\sigma_j)$ is nonincreasing, the step that opens a new epoch was computed with tolerance at most $\sigma_0$, and the tolerances used within any epoch are bounded by $\sigma_0$ and have square sum at most $\bar\Sigma$.

\subsection{Linear convergence}
\begin{theorem}\label{thm:acc-restart}
    Let \Cref{ass:acc-error-bound} hold and let the hypotheses of \Cref{thm:acc-inexact} hold with a nonincreasing tolerance sequence.
    Define
    \begin{align}
        C_\theta = \frac{2 + \sqrt{2m_\theta}\, e^{3c_\theta\bar\Sigma/2}}{\min\{1, \theta\}},
        \qquad
        J = \max\left\{2,\ \left\lceil \frac{C_\theta\, \eta\, (1 + \sigma_0)}{\rho}\right\rceil\right\}.
    \end{align}
    Then every epoch performs at most $J$ linear solves, the anchor steps satisfy $\|\bar s_{e+1}\|_M \le \rho\|\bar s_e\|_M$ for consecutive epochs $e$ and $e+1$, and the step produced by the $k$-th linear solve overall satisfies
    \begin{align}
        \|s^{(k)}\|_M \le C_\theta\,\eta\,(1 + \sigma_0)\; \rho^{\lfloor (k-1)/J \rfloor}\; \|\bar s_0\|_M, \qquad k \ge 1, \label{eq:acc-linear-rate}
    \end{align}
    where $\bar s_0$ is the step at the starting point.
\end{theorem}
\begin{proof}
    Consider an epoch with anchor $\bar z$ and opening step $\bar s = s^{(1)}$, and let $\bar D = \dist_M(\bar z, \mathcal{Z}^\star)$.
    Within the epoch the iterates follow the anchored method started at $\bar z$ with tolerances bounded by $\sigma_0$ and square-summable to at most $\bar\Sigma$, so \Cref{thm:acc-inexact} applies.
    Since $\theta j + 1 \ge \min\{1,\theta\}(j+1)$, its conclusion can be weakened to
    \begin{align}
        \|s^{(j+1)}\|_M \le \frac{C_\theta\,\bar D}{j+1} \qquad \text{for all } j \ge 0. \label{eq:acc-epoch-bound}
    \end{align}
    By \Cref{ass:acc-error-bound} and \Cref{eq:acc-exact-vs-inexact}, $\bar D \le \eta\|s^\star(\bar z)\|_M \le \eta(1+\sigma_0)\|s^{(1)}\|_M$.
    Therefore $\|s^{(j+1)}\|_M \le \rho\|s^{(1)}\|_M$ as soon as $j + 1 \ge C_\theta\eta(1+\sigma_0)/\rho$, so the restart test succeeds at some $j \le J - 1$, and the epoch performs at most $J$ linear solves (the solve that opens the epoch is charged to the previous one).
    The relation $\|\bar s_{e+1}\|_M \le \rho\|\bar s_e\|_M$ is the restart test itself.

    For \Cref{eq:acc-linear-rate}, note that the first epoch uses one solve for $\bar s_0$ and at most $J$ more, and every later epoch uses at most $J$.
    Hence the $k$-th solve belongs to an epoch $e \ge \lfloor (k-1)/J\rfloor$, whose anchor step satisfies $\|\bar s_e\|_M \le \rho^{e}\|\bar s_0\|_M$.
    Within that epoch, \Cref{eq:acc-epoch-bound} with $j \ge 0$ and the bound on $\bar D$ give $\|s^{(k)}\|_M \le C_\theta\bar D \le C_\theta\eta(1+\sigma_0)\|\bar s_e\|_M$.
\end{proof}

\begin{remark}[Comparison with the base method]
    Under \Cref{ass:acc-error-bound} the base method also converges linearly, but with a worse dependence on $\eta$.
    With exact solves, the draft's one-step inequality reads $\|z^{(k+1)} - z^\star\|_M^2 \le \|z^{(k)} - z^\star\|_M^2 - \theta(2-\theta)\|s^\star(z^{(k)})\|_M^2$ for every $z^\star \in \mathcal{Z}^\star$; choosing $z^\star$ as the projection of $z^{(k)}$ and applying the error bound gives
    \begin{align}
        \dist_M(z^{(k+1)}, \mathcal{Z}^\star)^2 \le \left(1 - \frac{\theta(2-\theta)}{\eta^2}\right)\dist_M(z^{(k)}, \mathcal{Z}^\star)^2 ,
    \end{align}
    so reaching accuracy $\epsilon$ takes on the order of $\eta^2\log(1/\epsilon)$ outer iterations.
    By \Cref{thm:acc-restart} the restarted anchored method takes on the order of $J\log(1/\epsilon)/\log(1/\rho) = O(\eta\log(1/\epsilon))$ outer iterations, which is the same improvement from $\eta^2$ to $\eta$ obtained by restarted Halpern methods for linear programming \citep{lu2024restarted}.
\end{remark}

\begin{remark}[Numerical sanity check]
    \mynote{TODO: turn into an experiment under \texttt{code/experiments/} with results saved as CSV.}
    On a random standard-form LP with $m = 60$ and $n = 200$, conjugate gradients on the Schur complement, $\theta = 1.5$, and constant $\sigma = 0.2$, the residual $\|s^{(k)}\|_M$ after $1000$ outer iterations was about $2.5\cdot 10^{-3}$ for the base method, $7\cdot 10^{-2}$ for the anchored method without restarts, and $8\cdot 10^{-6}$ for the restarted method with $\rho = 1/2$, which also used fewer inner iterations in total.
    This illustrates both halves of \Cref{sec:acc-restart}: without restarts the anchor is a liability on a problem where the base method converges linearly, and with restarts it is an asset.
\end{remark}

\section{Summary}
The anchored update \Cref{eq:anchored-update} accelerates the inexact primal-dual Douglas--Rachford method from $O(1/\sqrt{k})$ to $O(1/k)$ in the fixed-point residual and in both KKT residuals (\Cref{thm:acc-inexact,cor:acc-kkt}), at the price of a tolerance sequence with $\sum_k\sigma_k^2 < \infty$, which costs $O(\log k)$ inner iterations per outer iteration.
With exact solves the bound $2D/(\theta k+2)$ of \Cref{thm:acc-exact} is tight.
Restarting the anchor when the residual halves gives linear convergence under an error bound, with an iteration count linear rather than quadratic in the error-bound constant (\Cref{thm:acc-restart}).
Two questions remain open: whether a constant tolerance also gives $O(1/k)$ (\Cref{rem:acc-fixed-sigma}), and the transfer of the total inner-work bound of the draft to the anchored update (\Cref{rem:acc-inner-work}).

\bibliography{bibliography}

\end{document}
